\documentclass[11pt]{ctexart}

% ===================== 宏包 =====================
\usepackage{amsmath,amssymb,amsfonts}
\usepackage{geometry}
\geometry{a4paper,margin=2.4cm}
\usepackage{booktabs}
\usepackage{array}
\usepackage{enumitem}
\usepackage{xcolor}
\usepackage{tcolorbox}
\tcbuselibrary{skins,breakable}
\usepackage{listings}
\usepackage{hyperref}
\hypersetup{colorlinks=true,linkcolor=blue!55!black,urlcolor=teal,
            citecolor=blue!55!black,bookmarksopen=true}
\usepackage{fancyhdr}

% ===================== 颜色 =====================

\definecolor{accent}{RGB}{38,90,160}
\definecolor{warnbg}{RGB}{255,247,237}
\definecolor{warnrule}{RGB}{200,90,30}
\definecolor{notebg}{RGB}{238,246,255}
\definecolor{noterule}{RGB}{38,90,160}
\definecolor{codebg}{RGB}{246,248,250}
\definecolor{codekw}{RGB}{170,30,120}
\definecolor{codecm}{RGB}{100,120,100}

% ===================== 代码样式 =====================
\lstset{
  basicstyle=\ttfamily\small,
  backgroundcolor=\color{codebg},
  keywordstyle=\color{codekw}\bfseries,
  commentstyle=\color{codecm}\itshape,
  breaklines=true,
  frame=single,
  rulecolor=\color{black!15},
  framesep=5pt,
  xleftmargin=8pt,
  xrightmargin=4pt,
  showstringspaces=false,
  columns=fullflexible,
  keepspaces=true,
}

% ===================== 自定义盒子 =====================
\newtcolorbox{warnbox}[1]{
  enhanced,breakable,colback=warnbg,colframe=warnrule,
  boxrule=0.6pt,arc=2pt,left=8pt,right=8pt,top=6pt,bottom=6pt,
  fonttitle=\bfseries\color{warnrule},title={#1}}
\newtcolorbox{notebox}[1]{
  enhanced,breakable,colback=notebg,colframe=noterule,
  boxrule=0.6pt,arc=2pt,left=8pt,right=8pt,top=6pt,bottom=6pt,
  fonttitle=\bfseries\color{noterule},title={#1}}

% ===================== 页眉页脚 =====================
\pagestyle{fancy}
\fancyhf{}
\lhead{\small\color{accent}ARAIM-Aware EGO-Planner 集成手册}
\rhead{\small\thepage}
\renewcommand{\headrulewidth}{0.4pt}

\setlist{leftmargin=1.6em,itemsep=2pt,topsep=3pt}

% ===================== 标题 =====================
\title{\textbf{\color{accent}ARAIM 完整性感知的 EGO-Planner 集成手册}\\[4pt]
       \large 注入点设计、后端代价与运行时安全闭环（最终版）}
\author{规划与集成 / Integrity-Aware Planning}
\date{\today}

\begin{document}
\maketitle
\thispagestyle{fancy}

\begin{abstract}
\noindent 本手册定义将 ARAIM 完整性信息（$PL$、$AL$、$IM=AL-PL$）注入 EGO-Planner
的完整方案。核心原则：\textbf{安全陈述 $PL_{pred}<AL$ 是硬判定，只能由运行时
监督器（P5）裁决}；后端代价 P1--P4 仅提供\emph{偏好}，用于把轨迹引向定位条件更好的区域。
为保持 EGO 的无-ESDF 优势，含 $AL$ 的违例项不放入后端，而归 P5。本版相对前稿的主要变更：
新增并\textbf{前置 P5}、后端代价拆分修正（margin 项移出）、SafetyGrid 改为
raw/derived 分层、后端代价正式采用沿轨迹采样、时间维度与非光滑处理独立成节。
\end{abstract}

\tableofcontents
\bigskip

% =========================================================
\section{总体架构与原则}
% =========================================================

\subsection{两类语义的严格区分}
整个设计的正确性建立在区分以下两个概念之上：

\begin{center}\renewcommand{\arraystretch}{1.35}
\begin{tabular}{@{}p{4.3cm}p{4.8cm}p{4.6cm}@{}}
\toprule
 & \textbf{后端代价（preference）} & \textbf{安全条件（safety）}\\
\midrule
形式 & $\lambda_{pref}\psi(PL_{pred})$ & $PL_{pred}(p(t),t)<AL(p(t))$\\
作用 & 偏好定位质量更高的区域 & 判定轨迹是否完整性安全\\
性质 & localization-quality preference & integrity safety condition\\
归属 & 后端优化器（P1--P4） & 运行时监督器（\textbf{P5}）\\
\bottomrule
\end{tabular}\end{center}

\begin{warnbox}{切勿混淆}
不要把 ``minimize $PL_{pred}$'' 说成``满足 safety''。最小化 $PL$ 只能让轨迹
\emph{偏向高完整性区域}；真正的 safety statement 必须来自 P5 对
$PL_{pred}<AL$ 的逐采样点硬判定。
\end{warnbox}

\subsection{注入点总览与实现优先级}
\textbf{P5 排在 P1 之前}：硬安全闭环是``系统不会撞 $PL$ 墙''的底线，应在任何
偏好优化之前就位。有了 P5，才能放心地逐步上 P1--P4 这些会改变轨迹形状的偏好项。

\begin{center}\renewcommand{\arraystretch}{1.3}
\begin{tabular}{@{}clll@{}}
\toprule
顺序 & 注入点 & 作用 & 性质\\
\midrule
Step 0 & P0 风险栅格 (raw+derived) & 前置基础设施 & 数据层\\
Step 1 & \textbf{P5 Integrity Supervisor} & $PL<AL$ 硬闭环 & \textbf{安全底线（优先）}\\
Step 2 & P1 后端 integrity 代价 & 走廊内贴安全侧 & EGO-native\\
Step 3 & P2 候选排序 & 已有候选选优 & 性价比最高\\
Step 4 & P3 全局走廊偏置 & 任务级走廊 & 近乎免费\\
Step 5 & P4 risk-aware A$^\ast$ 初值 & 生成新走廊 & 仅按需\\
\bottomrule
\end{tabular}\end{center}

% =========================================================
\section{P0 —— 风险栅格 SafetyGrid（raw/derived 分层）}
% =========================================================

\subsection{设计动机}
SafetyGrid 是所有完整性注入点的数据底座。它把 predictor 的输出空间化为可被
planner 高速查询的栅格场。

\paragraph{分层存储（不要只存 $c_{PI}$）。}
\begin{itemize}
\item \textbf{Raw layers}（不依赖 planner 策略，可直接缓存/序列化）：
  \[
  hpl_{adv},\; vpl_{adv},\; pl_{scalar},\; source\_flags,\;
  valid/available/fallback,\; fallback\_reason,\; age/generation .
  \]
\item \textbf{Derived layers}（依赖 planner 策略，可\emph{惰性计算 + 缓存}）：
  \[
  AL,\; IM,\; c_{PI},\; c_{unknown},\; risk\_band,\; \nabla_p c_{PI} .
  \]
\end{itemize}

\begin{notebox}{为什么不能只存 $c_{PI}$}
$c_{PI}$ 依赖 $AL$、margin $m$、权重、unknown policy、是否启用 PL-only preference
等\emph{planner 策略}。若栅格只存 $c_{PI}$，则后续调参或换策略必须重建整张栅格。
分层后，raw 层稳定复用，derived 层随策略惰性重算，二者解耦。
\end{notebox}

\subsection{查询结果的类型拆分}
为避免 prediction、grid metadata、planner cost 三者混在一个结构里，查询接口拆为：
\begin{lstlisting}[language=C++]
struct PredictorQueryResult { /* hpl, vpl, source_flags, ... */ };
struct GridQueryResult      { /* valid, age, generation, fallback_reason */ };
struct PlannerRiskResult    { /* AL, IM, c_PI, grad_c_PI, risk_band */ };
\end{lstlisting}

% =========================================================
\section{P5 —— Runtime Integrity Supervisor（硬安全闭环，唯一裁决者）}
% =========================================================

\begin{warnbox}{为什么必须独立存在}
P1--P4 都是\emph{偏好}：让轨迹倾向定位更好的区域。它们无法回答
``如果未来轨迹已不满足 $PL_{pred}<AL$ 怎么办''。安全陈述
$PL_{pred}(p(t),t)<AL(p(t))$ 的\textbf{硬判定权必须且只能}归 P5。
没有 P5，系统就只是``优化器尝试偏好低风险路线''，却缺少安全闭环。
\end{warnbox}

\paragraph{位置。}\texttt{ego\_replan\_fsm.cpp} 的轨迹安全检查 / 状态机回调；
与 EGO 原有 \texttt{checkTrajCollision()} 并列。

\paragraph{输入。}当前 ARAIM 监测量 $PL_{mon}$、当前 $AL$、对\emph{已选轨迹}
逐采样点的 $PL_{pred}(p(t_r),t_r)$ 与 $AL(p(t_r))$（其中 $IM(t_r)=AL(t_r)-PL_{pred}(t_r)$）。

\paragraph{裁决逻辑。}
\begin{lstlisting}
if (PL_mon > AL_current)            -> EMERGENCY (stop / hover);
if (min_r IM(t_r) < 0)             -> trigger replan or slow down;
if (safety_grid stale/unknown ||
    predictor unavailable)         -> conservative fallback + unknown penalty;
\end{lstlisting}

\paragraph{改动清单。}
\begin{center}\renewcommand{\arraystretch}{1.3}
\begin{tabular}{@{}p{5.2cm}p{8.4cm}@{}}
\toprule 文件 / 函数 & 改法 \\
\midrule
\texttt{ego\_replan\_fsm.h/.cpp} & 新增 \texttt{checkTrajIntegrity()}；状态机加
  \texttt{EMERGENCY\_HOVER} / \texttt{SLOW\_DOWN} 转移\\
\texttt{PlannerAdapter} & 提供 $IM=AL-PL_{pred}$ 的逐采样点查询\\
\bottomrule\end{tabular}\end{center}

% =========================================================
\section{P1 —— 后端 integrity 代价（EGO-native 偏好项）}
% =========================================================

\subsection{代价定义（正式方案：沿轨迹采样）}
\begin{equation}
J_{PI}^{traj}(Q)=\sum_{r=1}^{N_s} c_{PI}\big(p(t_r),t_r\big)\,\Delta t_r,
\qquad p(t_r)=\sum_i B_i^k(t_r)\,Q_i,
\end{equation}
对控制点 $Q_i$ 的梯度（map-gradient $\times$ B-spline basis 链式法则）：
\begin{equation}
\frac{\partial J_{PI}^{traj}}{\partial Q_i}
 =\sum_{r=1}^{N_s}\nabla_p c_{PI}\big(p(t_r),t_r\big)\,B_i^k(t_r)\,\Delta t_r .
\end{equation}

\begin{notebox}{控制点版 vs 采样版（升级有明确触发条件）}
\textbf{默认用控制点版} $J_{PI}=\sum_i c_{PI}(Q_i)$：它与 EGO 自身
``基于凸包的控制点评估''哲学一致，对一个\emph{软偏好项}通常够用、且更廉价。\\
\textbf{升级到采样版的触发条件}：P5 观测到``控制点低风险但曲线中段 $PL_{pred}\ge AL$''。\\
\textbf{数学原因}：占据约束配合凸包 + 膨胀半径，故控制点评估成立；而 risk field
\emph{非凸、无膨胀语义}，控制点低风险不能保证 B-spline 曲线内部也低风险。
\end{notebox}

\subsection{后端代价的正确拆分（关键概念）}
\begin{equation}
J_I=\underbrace{\lambda_{pref}\,\psi(PL_{pred})}_{\text{定位质量偏好}}
   +\underbrace{\lambda_{unk}\,J_{unknown}}_{\text{stale / unknown 处理}} .
\end{equation}

\begin{warnbox}{后端不放含 $AL$ 的 margin 项}
$AL=\min(\gamma_H d_{obs},\gamma_V d_{vertical})$ 依赖障碍距离。把
$[PL_{pred}-AL+m]_+^2$ 放进后端，等于把 ESDF 依赖\emph{偷偷请回}后端，
破坏 EGO 的无-ESDF 优势。因此：
\begin{itemize}
\item 真正的违例项 $[PL_{pred}-AL+m]_+$ 归 \textbf{P5}（评估器本就逐点判 $PL<AL$）；
\item 障碍 clearance 仍交给 EGO 原有的 rebound collision cost；
\item 若工程上确需后端 margin，则 $AL_{proxy}$ 必须是\emph{不查实时 ESDF}的廉价代理
  （常数 $AL$，或粗 grid 预存的 derived layer），\textbf{不得}实时查 ESDF。
\end{itemize}
\end{warnbox}

\subsection{与碰撞代价的优先级}
$\lambda_{PI}$ 采用退火 + 归一化 + 饱和，并通过 gradient clipping 保证
integrity 梯度\textbf{不压过} collision 梯度——碰撞安全永远优先于完整性偏好。

% =========================================================
\section{Integrity 代价的非光滑处理（L-BFGS 前提）}
% =========================================================

\subsection{非光滑来源}
\begin{itemize}
\item GNSS 可见卫星集（visible satellite set）切换；
\item LiDAR primitive 集切换；
\item fused source fallback 切换；
\item $\max(\text{GNSS},\text{LiDAR},\text{fused})$ 或 $\max$ over advisory 退化模式；
\item unknown / stale flag 翻转；
\item hinge loss $[\cdot]_+$。
\end{itemize}
直接插值 $c_{PI}$ 会产生跳变梯度，毁掉 L-BFGS 的曲率（拟牛顿）估计。

\subsection{必做的平滑措施}
\begin{enumerate}
\item $c_{PI}$ 做\textbf{饱和}（saturation），避免极大 $PL$ 产生爆炸梯度；
\item hinge 用 \textbf{smooth-hinge / pseudo-Huber}；
\item $\max$ 用 \textbf{log-sum-exp / softmax} 近似（至少在 planner cost 层平滑）；
\item \textbf{切换边界直接标 transition / unknown flag}，在该处走 unknown penalty，
  而\emph{非}去追逐不连续边的伪梯度；
\item \textbf{gradient clipping}，保证 integrity 梯度不会压过 collision 梯度。
\end{enumerate}

% =========================================================
\section{P2 —— Candidate ranking over existing initial guesses（候选排序）}
% =========================================================

\paragraph{能力。}从已有候选中选择更低 risk 的 basin。候选集通常包括：
previous-trajectory 续接、global-reference 采样、random-polynomial fallback。

\paragraph{限制。}\textbf{不保证}发现新的 homotopy class。只有当候选生成器本身
吐出多条不同走廊候选时，P2 才真正具备``跨走廊选择''能力。

\begin{warnbox}{不要过度声称``跨走廊''}
若候选集仅为上述三类，P2 本质是 candidate scoring / selection——它能选出
候选集中更安全的一条，但不能产生新的同伦类。
\end{warnbox}

\paragraph{升级条件。}若候选集没有覆盖安全走廊，启用 \textbf{P4} risk-aware A$^\ast$
重新生成走廊候选。

% =========================================================
\section{P3 —— 全局参考偏置（任务级走廊）}
% =========================================================
在全局参考路径采样阶段，按 SafetyGrid 的 $risk\_band$ 对参考点施加偏置，
使局部规划在任务级别就倾向高完整性走廊。该项近乎免费（复用既有全局参考管线），
仅改变参考采样分布，不引入新优化变量。

% =========================================================
\section{P4 —— Risk-aware A$^\ast$ 初值生成（按需）}
% =========================================================
当 P2 判定现有候选集未覆盖安全走廊时启用。在 A$^\ast$ 的边代价中加入
$c_{PI}$ 项，使搜索能\emph{生成}穿越高完整性区域的新初值，从而产生新的同伦类
供后端优化。该项成本最高，仅按需触发。

% =========================================================
\section{时间维度（predictor 语义 $f_{pred}(p,t+\tau)$）}
% =========================================================

\begin{notebox}{物理判断：为什么准静态在短 horizon 内成立}
纯 GNSS 几何（可见集 / DOP）随时间\textbf{分钟级}变化，在 $\le 3$\,s 的规划
horizon 内几乎不变。因此``\textbf{空间变化 + 时间准静态}''是主导近似。
真正快变且位置相关的效应（建筑遮挡导致的可见集突变、近场多径、LiDAR 退化）
主要已由\emph{空间}维度捕获——因为飞机移动到未来位置 $p$ 时，空间查询已反映该变化。
故 Stage 1 并非``粗糙凑合'',对 GNSS 几何分量是有充分依据的良好近似；多 horizon
层主要为城市峡谷下可见集会切换的场景留升级口。
\end{notebox}

\paragraph{实现分级。}
\begin{center}\renewcommand{\arraystretch}{1.3}
\begin{tabular}{@{}ll@{}}
\toprule
Stage 1 & quasi-static SafetyGrid，$query\_time=$ snapshot.stamp，$horizon=0$\\
Stage 2 & multi-horizon 层 $\tau=0,1,2,3$\,s，optimizer 按采样时刻在层间插值\\
Stage 3 & 仅对\emph{最终选定}轨迹做 exact time-aware 校验（在 P5 内执行）\\
\bottomrule
\end{tabular}\end{center}

\begin{warnbox}{硬约定}
$query\_time\_s$ 与 $horizon\_s$ \textbf{必须是 finite}，\emph{不得用 NaN 表示
``当前时间''}。该约定从 predictor 继承到 SafetyGrid 与 PlannerAdapter。
\end{warnbox}

% =========================================================
\section{接口与职责边界小结}
% =========================================================
\begin{center}\renewcommand{\arraystretch}{1.3}
\begin{tabular}{@{}lll@{}}
\toprule
模块 & 持有 & 不持有\\
\midrule
Predictor & $hpl,vpl,source\_flags$（raw 预测） & planner 策略 / cost\\
SafetyGrid & raw layers + 缓存的 derived layers & 安全裁决权\\
PlannerAdapter & $AL,IM,c_{PI},\nabla c_{PI}$（derived） & 原始预测语义\\
后端优化器 (P1--P4) & 偏好代价 $J_I$（无 $AL$ 项） & $PL<AL$ 硬判定\\
\textbf{P5 Supervisor} & \textbf{$PL<AL$ 唯一硬判定权} & 形状级优化\\
\bottomrule
\end{tabular}\end{center}

\bigskip
\begin{notebox}{一句话总结}
后端管``往定位好的一侧贴''（preference，无 ESDF 依赖）；
P5 管``是否真的满足 $PL_{pred}<AL$''（safety，唯一裁决）。
两者职责不交叉，是整套设计正确性的基石。
\end{notebox}

\end{document}